本文以 OpenAI 于 2026 年 9 月 23 日、24 日发布的两篇挂谷猜选手稿为主要对象：一篇声称证明三维挂谷极大函数猜想，把 Wang--Zahl 定理留下的密度幂缺口 $\lambda^{K(\varepsilon)}$ 补到尖锐的 $\lambda^{3}$；另一篇声称证明四维挂谷集的 Hausdorff 维数等于 $4$，且不设可测性、紧性与粘性假设。两篇手稿均未经同行评审，也未纳入 Lean 形式化清单，仓库自述未形式化的结果可能存在问题，因此本文的第一项工作是事实核验：整理发布时间线、稿件参数、依赖链与验证状态，并指出四维结果的关键输入取自三维姊妹篇的一条加权板条引理，其有效性无法与三维篇的核验状态切割开来。第二项工作是技术拆解：从任意挂谷集到加权模型的归约、以“地平线”为载体的时间代价记账、带质量约束的二次图卡、极值临界路径与三窄度区间的分类讨论，逐层展开四维证明的机理，并说明它与三维范式（粘性、平面性、颗粒性）在对象层面的区别。第三项工作是技术谱系分析：把核心技术分为直系继承、门内改造、跨门派借入与谱系内首次系统使用四类，逐条给出证据与置信度。对照 Wang--Zahl 与 Guth--Wang--Zahl 原文可以确认，手稿使用的熵账本、概率耦合、实代数几何的定量纪律与乘法卷积的 Fourier 衰减估计，在传统挂谷工具箱中并无先例。最后给出可复现的核验路线图，以及验证、数学推进、方法回灌、工程流程四个方向的下一步建议。全文对手稿声称与已证定理严格区分。